% GENERATED FILE. Edit canonical lessons, then run npm run generate:books.
% canonical-source-hash: fnv1a64:6f9e16a4521f6126
% canonical-lessons: FA-C382-prozhe, FA-C382-modiriyyat-kardan, FA-C382-barname-rizi-kardan, FA-C382-sazman, FA-C382-herfei

\chapter{Projects and planning}
\label{ch:fa-projects-and-planning}
\hlchaptermodality{382}

\begin{chapteropening}
\textbf{By the end of this chapter:} \emph{I can talk about a project, say who manages it, how we planned it and which organisation is behind it.}

Complete the last lesson of chapter 382: I can talk about a project, say who manages it, how we planned it and which organisation is behind it.
\end{chapteropening}

\section[\texorpdfstring{\fa{پروژه} (prozhe)}{پروژه (prozhe)}]{\fa{پروژه} (prozhe) — a project}

\label{lesson:FA-C382-prozhe}



\begin{quote}
Before the new one: say the Persian for a specialist, then the Persian for a duty.
\end{quote}

\subsection*{You'll want to know: \fa{پروژه}}
\textbf{\fa{پروژه}} (\emph{prozhe}) — \textquotedblleft{}a project\textquotedblright{}.

\begin{quote}
\fa{این} \fa{پروژه} \fa{شش} \fa{ماه} \fa{طول} \fa{میکشد}.
\end{quote}

\begin{quote}
\emph{in prozhe shesh mâh tul mikeshad.}
\end{quote}

\begin{quote}
This project takes six months.
\end{quote}

A loan from French. Plural \emph{prozhehâ}. \fa{ژ} sounds like the s in English pleasure.

\begin{cousinweb}
\textbf{\fa{پروژه}} was borrowed from French \emph{projet} \textquotedblleft{}project\textquotedblright{} (EtymDB), from Latin \emph{prōiectus} \textquotedblleft{}thrown forward\textquotedblright{}, built on \emph{iacere} \textquotedblleft{}to throw\textquotedblright{} (Lewis \& Short): a plan thrown ahead of you.
\end{cousinweb}

\begin{grammarlens}[title={What you've built}]
One more word for telling, explaining and describing at B1.
\end{grammarlens}

\subsection*{Your turn}
\begin{itemize}
\raggedright
  \item \emph{Say it:} \emph{prozhe}
  \item \emph{Say it:} \emph{prozhe}, once more
  \item \emph{Say it:} a sentence of your own with \emph{prozhe}
  \item \emph{Recall:} say the Persian for a specialist, then the Persian for a duty, then say \emph{prozhe} again
\end{itemize}

\begin{tcolorbox}[breakable,colback=teal!4,colframe=teal!35!black,title={Before you move on}]
What does \fa{پروژه} mean? (\textquotedblleft{}A project\textquotedblright{}.) Say it once more.
\end{tcolorbox}
\section[\texorpdfstring{\fa{مدیریت} \fa{کردن} (modiriyat kardan)}{مدیریت کردن (modiriyat kardan)}]{\fa{مدیریت} \fa{کردن} (modiriyat kardan) — to manage, to run}

\label{lesson:FA-C382-modiriyyat-kardan}



\begin{quote}
Before the new one: say the Persian for a duty, then the Persian for a project.
\end{quote}

\subsection*{You'll want to know: \fa{مدیریت} \fa{کردن}}
\textbf{\fa{مدیریت} \fa{کردن}} (\emph{modiriyat kardan}) — \textquotedblleft{}to manage, to run\textquotedblright{}.

\begin{quote}
\fa{مادرم} \fa{یک} \fa{رستوران} \fa{کوچک} \fa{را} \fa{مدیریت} \fa{میکند}.
\end{quote}

\begin{quote}
\emph{mâdaram yek resturân-e kuchek râ modiriyat mikonad.}
\end{quote}

\begin{quote}
My mother runs a small restaurant.
\end{quote}

\emph{Modiriyyat} \textquotedblleft{}management\textquotedblright{}, built on \emph{modir} \textquotedblleft{}manager\textquotedblright{}, + \emph{kardan}.

\begin{cousinweb}
\textbf{\fa{مدیریت}} is built on \emph{modir} \textquotedblleft{}manager\textquotedblright{}, borrowed from Arabic \emph{mudīr}, with the abstract ending \emph{-iyat}; \emph{kardan} \textquotedblleft{}to do\textquotedblright{} is the inherited Persian verb.
\end{cousinweb}

\begin{grammarlens}[title={What you've built}]
One more word for telling, explaining and describing at B1.
\end{grammarlens}

\subsection*{Your turn}
\begin{itemize}
\raggedright
  \item \emph{Say it:} \emph{modiriyat kardan}
  \item \emph{Say it:} \emph{modiriyat kardan}, once more
  \item \emph{Say it:} a sentence of your own with \emph{modiriyat kardan}
  \item \emph{Recall:} say the Persian for a duty, then the Persian for a project, then say \emph{modiriyat kardan} again
\end{itemize}

\begin{tcolorbox}[breakable,colback=teal!4,colframe=teal!35!black,title={Before you move on}]
What does \fa{مدیریت} \fa{کردن} mean? (\textquotedblleft{}To manage, to run\textquotedblright{}.) Say it once more.
\end{tcolorbox}
\section[\texorpdfstring{\fa{برنامه} \fa{ریزی} \fa{کردن} (barnâme rizi kardan)}{برنامه ریزی کردن (barnâme rizi kardan)}]{\fa{برنامه} \fa{ریزی} \fa{کردن} (barnâme rizi kardan) — to plan, to schedule}

\label{lesson:FA-C382-barname-rizi-kardan}



\begin{quote}
Before the new one: say the Persian for a project, then the Persian for to manage.
\end{quote}

\subsection*{You'll want to know: \fa{برنامه} \fa{ریزی} \fa{کردن}}
\textbf{\fa{برنامه} \fa{ریزی} \fa{کردن}} (\emph{barnâme rizi kardan}) — \textquotedblleft{}to plan, to schedule\textquotedblright{}.

\begin{quote}
\fa{ما} \fa{برای} \fa{تعطیلات} \fa{تابستان} \fa{برنامه} \fa{ریزی} \fa{کردیم}.
\end{quote}

\begin{quote}
\emph{mâ barâye ta'tilât-e tâbestân barnâme rizi kardim.}
\end{quote}

\begin{quote}
We made plans for the summer holidays.
\end{quote}

\emph{Barnâme} \textquotedblleft{}programme\textquotedblright{} + \emph{rizi} \textquotedblleft{}pouring\textquotedblright{}, from \emph{rikhtan} \textquotedblleft{}to pour\textquotedblright{}, + \emph{kardan}.

\begin{cousinweb}
\textbf{\fa{برنامه} \fa{ریزی} \fa{کردن}} is native Persian, built from \emph{barnâme} \textquotedblleft{}programme\textquotedblright{}, which EtymDB traces to Middle Persian \emph{barnāmag}, + \emph{rizi}, built on \emph{riz-}, the present stem of \emph{rikhtan} \textquotedblleft{}to pour\textquotedblright{}, + \emph{kardan}. To plan is to pour out a programme.
\end{cousinweb}

\begin{grammarlens}[title={What you've built}]
One more word for telling, explaining and describing at B1.
\end{grammarlens}

\subsection*{Your turn}
\begin{itemize}
\raggedright
  \item \emph{Say it:} \emph{barnâme rizi kardan}
  \item \emph{Say it:} \emph{barnâme rizi kardan}, once more
  \item \emph{Say it:} a sentence of your own with \emph{barnâme rizi kardan}
  \item \emph{Recall:} say the Persian for a project, then the Persian for to manage, then say \emph{barnâme rizi kardan} again
\end{itemize}

\begin{tcolorbox}[breakable,colback=teal!4,colframe=teal!35!black,title={Before you move on}]
What does \fa{برنامه} \fa{ریزی} \fa{کردن} mean? (\textquotedblleft{}To plan, to schedule\textquotedblright{}.) Say it once more.
\end{tcolorbox}
\section[\texorpdfstring{\fa{سازمان} (sâzmân)}{سازمان (sâzmân)}]{\fa{سازمان} (sâzmân) — an organisation}

\label{lesson:FA-C382-sazman}



\begin{quote}
Before the new one: say the Persian for to manage, then the Persian for to plan.
\end{quote}

\subsection*{You'll want to know: \fa{سازمان}}
\textbf{\fa{سازمان}} (\emph{sâzmân}) — \textquotedblleft{}an organisation\textquotedblright{}.

\begin{quote}
\fa{او} \fa{در} \fa{یک} \fa{سازمان} \fa{بین} \fa{المللی} \fa{کار} \fa{میکند}.
\end{quote}

\begin{quote}
\emph{u dar yek sâzmân-e beyn ol-melali kâr mikonad.}
\end{quote}

\begin{quote}
She works for an international organisation.
\end{quote}

From \emph{sâz-}, the present stem of \emph{sâkhtan} \textquotedblleft{}to make\textquotedblright{}, + \emph{-mân}. Plural \emph{sâzmânhâ}.

\begin{cousinweb}
\textbf{\fa{سازمان}} is native Persian, built on \emph{sâz-}, the present stem of \emph{sâkhtan} \textquotedblleft{}to make, to build\textquotedblright{}, with the suffix \emph{-mân}: an organisation is something built and arranged.
\end{cousinweb}

\begin{grammarlens}[title={What you've built}]
One more word for telling, explaining and describing at B1.
\end{grammarlens}

\subsection*{Your turn}
\begin{itemize}
\raggedright
  \item \emph{Say it:} \emph{sâzmân}
  \item \emph{Say it:} \emph{sâzmân}, once more
  \item \emph{Say it:} a sentence of your own with \emph{sâzmân}
  \item \emph{Recall:} say the Persian for to manage, then the Persian for to plan, then say \emph{sâzmân} again
\end{itemize}

\begin{tcolorbox}[breakable,colback=teal!4,colframe=teal!35!black,title={Before you move on}]
What does \fa{سازمان} mean? (\textquotedblleft{}An organisation\textquotedblright{}.) Say it once more.
\end{tcolorbox}
\section[\texorpdfstring{\fa{حرفه} \fa{ای} (herfe'i)}{حرفه ای (herfe'i)}]{\fa{حرفه} \fa{ای} (herfe'i) — professional}

\label{lesson:FA-C382-herfei}



\begin{quote}
Before the new one: say the Persian for to plan, then the Persian for an organisation.
\end{quote}

\subsection*{You'll want to know: \fa{حرفه} \fa{ای}}
\textbf{\fa{حرفه} \fa{ای}} (\emph{herfe'i}) — \textquotedblleft{}professional\textquotedblright{}.

\begin{quote}
\fa{او} \fa{یک} \fa{راننده} \fa{حرفه} \fa{ای} \fa{است}.
\end{quote}

\begin{quote}
\emph{u yek rânande-ye herfe'i ast.}
\end{quote}

\begin{quote}
He is a professional driver.
\end{quote}

\emph{Herfe} \textquotedblleft{}trade, profession\textquotedblright{} + \emph{-i}, written apart. The opposite is \emph{âmâtor} or \emph{gheyr-e herfe'i}.

\begin{cousinweb}
\textbf{\fa{حرفه} \fa{ای}} is \emph{herfe}, from Arabic \emph{ḥirfa} \textquotedblleft{}occupation\textquotedblright{} (EtymDB), with the Persian adjective suffix \emph{-i}.
\end{cousinweb}

\begin{grammarlens}[title={What you've built}]
One more word for telling, explaining and describing at B1.
\end{grammarlens}

\subsection*{Your turn}
\begin{itemize}
\raggedright
  \item \emph{Say it:} \emph{herfe'i}
  \item \emph{Say it:} \emph{herfe'i}, once more
  \item \emph{Say it:} a sentence of your own with \emph{herfe'i}
  \item \emph{Recall:} say the Persian for to plan, then the Persian for an organisation, then say \emph{herfe'i} again
\end{itemize}

\begin{tcolorbox}[breakable,colback=teal!4,colframe=teal!35!black,title={Before you move on}]
What does \fa{حرفه} \fa{ای} mean? (\textquotedblleft{}Professional\textquotedblright{}.) Say it once more.
\end{tcolorbox}
